\documentclass[11pt]{article}
\usepackage{icmlko}
\begin{document}

\title{하루 만에 되돌렸다 --- 씨앗 스물이 모자랐다}
\author{월드모델 실험실}
\date{노트 507 · 2026년 8월 3일}
\maketitle

\begin{abstract}
노트 504 가 세계애니 \texttt{grp} 을 채택하면서 되돌림 조건을 미리 적었다
--- \emph{짝 씨앗 $40$ 으로 다시 재서 판 구간이 $0$ 을 물면 되돌린다.}

\textbf{발동했다.}

\begin{center}
\begin{tabular}{lrrl}
\hline
자 & 짝 씨앗 $20$ & 짝 씨앗 $40$ & $95\%$ 구간($40$) \\
\hline
판 & $+0.0011$ & $\mathbf{+0.0000}$ $22/40$ & $[-0.0004, +0.0005]$ \\
날짜 통제 판 & $+0.0004$ & $\mathbf{-0.0014}$ $8/40$ & $[-0.0019, -0.0009]$ \\
무리 안 판 & $+0.0004$ & $-0.0006$ $17/40$ & $[-0.0013, +0.0002]$ \\
KR 만화 & $+0.0084$ & $+0.0068$ $32/40$ & $[+0.0046, +0.0090]$ \\
비게임 앱 & $-0.0020$ & $\mathbf{-0.0032}$ $10/40$ & $[-0.0048, -0.0015]$ \\
\hline
\end{tabular}
\end{center}

씨앗 $20$ 에서 판 $+0.0011$ 이 구간 $[+0.0003, +0.0019]$ 로 $0$ 밖이었다.
씨앗 $40$ 에서 \textbf{정확히 $0$} 이다. 그리고 씨앗 $20$ 에서 ``$0$ 을
문다''던 날짜 통제 판과 앱이 \textbf{둘 다 확실한 음수}가 됐다.

\textbf{같은 후보 · 같은 자료 · 씨앗만 두 배다.} 되돌렸다 ---
\texttt{DOMDROP = \{"애니": ("trend\_",)\}}.
\end{abstract}

\begin{figure}[t]
\centering
\includegraphics[width=\textwidth]{figs/f1.pdf}
\caption{\textbf{가} --- 자 다섯, 씨앗 $20$ 대 $40$. \textbf{나} ---
판 하나만. \textbf{다} --- 두 채택의 크기.}
\end{figure}

\section{무엇을 배웠나}

노트 504 는 규율을 다 지켰다 --- 사전 등록을 실행 전에 썼고, 통과 기준을
자 다섯으로 두었고, 약하다는 것을 본문에 적었고, \textbf{되돌림 조건을
미리 박았다.} 그런데도 틀렸다.

틀린 곳은 하나다 --- \textbf{씨앗 $20$ 이 그 크기의 효과에 모자랐다.}
판 차 $+0.0011$ 은 짝 씨앗 잡음과 같은 자리에 있다.

\begin{quote}
\textbf{조항} --- \textbf{판 차가 $+0.002$ 아래인 후보는 짝 씨앗 $40$ 으로
재기 전에 채택하지 않는다.} 노트 498 의 애니 트렌드($+0.0045$)는 안
걸린다. 노트 $447 \cdot 449$ 의 계절 $\cdot$ 특기($+0.0041 \cdot
+0.0049$)도 안 걸린다. \textbf{이 문턱 아래로 내려간 채택은 세계애니
\texttt{grp} 하나뿐이었고 그것이 틀렸다.}
\end{quote}

\section{되돌림 조건이 일했다}

이 실험실이 되돌림 조건을 \emph{미리} 적기 시작한 것은 노트 432 부터다
(``되돌림 조항은 `유지'라 쓰지 말고 이전 상태를 이름으로 적는다''). 그
규약이 \textbf{처음으로 실제 채택을 하루 만에 잡았다.}

\begin{quote}
채택할 때 ``이것이 틀렸다면 무엇으로 알게 되나''를 같이 적으면, 틀렸을 때
\emph{다음 사이클에} 알게 된다. 안 적으면 \emph{영영} 모른다.
\end{quote}

\section{그래서 지금 상태}

\begin{center}
\begin{tabular}{lrl}
\hline
채택 & 판 차 & 상태 \\
\hline
계절 축(노트 447) & $+0.0041$ & 유지 \\
특기 칸(노트 449) & $+0.0049$ & 유지 \\
깊이 $12$(노트 457) & $-0.0032$ & 유지(집 밖으로 삼) \\
애니 트렌드(노트 498) & $+0.0045$ & \textbf{유지} \\
세계애니 \texttt{grp}(노트 504) & $+0.0011$ & \textbf{되돌림} \\
\hline
\end{tabular}
\end{center}

판은 $0.4854 \to \mathbf{0.4857}$(노트 498 상태)로 돌아간다.

\section*{남은 것}

\begin{itemize}
\item \textbf{노트 $498$ 애니 트렌드도 씨앗 $40$ 으로} --- $+0.0045$ 라
새 문턱에 안 걸리지만, 같은 실행에서 확인해 두면 값싸다.
\item \textbf{씨앗 $20$ 으로 판정한 과거 채택이 또 있나} --- 계절 $\cdot$
특기는 씨앗 $12$ 였다. 크기가 문턱 위($+0.0041 \cdot +0.0049$)라 규칙상
안 걸리지만, 그 규칙 자체가 오늘 만든 것이다.
\item \textbf{자료 쪽으로} --- 노트 $505 \cdot 506$ 이 모형 갈래를 닫고
행 늘리기를 가리켰다. 그쪽이 이 $\pm 0.001$ 싸움보다 스물다섯 배 크다.
\end{itemize}

\end{document}
